% Results section -- PhD4 (Results and Discussion specialist)
% All numbers traceable to CSVs in project path /usr/fafh.
% Source: model_comparison_national.csv, results_national_*.csv, data/imputation_selection_results.csv,
%         results_national_*_robust.csv, model_accuracy_summary.csv.

\section{Results}
\label{sec:results}

\subsection{Main Estimates of National Out-of-Home Consumption Coefficients}

We estimate national out-of-home consumption coefficients for 15 food categories and 10 years (2015--2024) using 12 prediction models. The coefficient for each category is the ratio of total (at-home plus out-of-home) consumption to at-home consumption; a value above one indicates that consumption away from home raises total consumption above the at-home level. All estimates in this subsection are taken from \texttt{model\_comparison\_national.csv} and the \texttt{results\_national\_*.csv} files (e.g., \texttt{results\_national\_lightgbm.csv}, \texttt{results\_national\_tabpfn.csv}, \texttt{results\_national\_fttransformer.csv}).

Across models and years, grain categories (rice \textit{daogu}, wheat \textit{xiaomai}, legumes \textit{doulei}, coarse grains \textit{zaliang}) show coefficients typically between about 1.05 and 1.35. For example, in 2020 LightGBM yields 1.128 for rice, 1.153 for wheat, 1.182 for legumes, and 1.153 for coarse grains (\texttt{model\_comparison\_national.csv}, rows for Year=2020, model=lightgbm). Vegetables (\textit{shucai}) are generally in the 1.05--1.19 range (e.g., LightGBM 2020: 1.069; TabPFN 2024: 1.187). Meat and animal products show larger multipliers: pork (\textit{zhurou}) often between 1.16 and 1.38, poultry (\textit{qinlei}) and beef (\textit{niurou}) in similar or higher ranges depending on model and year. Egg (\textit{danlei}), dairy (\textit{nailei}), fruits (\textit{guaguo}), cooking oil (\textit{youliao}), and sugar (\textit{tang}) coefficients are usually closer to one or modestly above (e.g., oil and sugar near 1.00--1.02 in many specifications).

The tree-based and neural/table models give broadly consistent orders of magnitude, but levels differ. LightGBM tends to produce somewhat lower coefficients for grains and meat in later years (e.g., 2024 rice 1.110, pork 1.168), whereas TabPFN and FT-Transformer often give higher values (e.g., TabPFN 2024 rice 1.233, pork 1.376; FT-Transformer 2024 rice 1.242, pork 1.228). Linear and Lasso specifications sit in the middle for many categories; some flexible models (e.g., XGBoost, ResNet, linear in a few categories) exhibit unstable or extreme point estimates in particular year--category cells and are therefore less reliable for policy interpretation. We emphasize LightGBM, TabPFN, and FT-Transformer for their combination of prediction accuracy and stability across categories: TabPFN attains the best prediction MAE, FT-Transformer is the best imputation model and ranks first on RMSE for the prediction task, and LightGBM provides stable point estimates consistent with the tabular-deep-learning models.

Table~\ref{tab:national_coef_summary} reports national out-of-home consumption coefficients for selected food categories and years from LightGBM and TabPFN. Every entry is taken from \texttt{model\_comparison\_national.csv} (columns \texttt{全国户外消费系数\_daogu}, \texttt{全国户外消费系数\_xiaomai}, \texttt{全国户外消费系数\_zhurou}, \texttt{全国户外消费系数\_shucai}).

\begin{table}[htbp]
\centering
\caption{National out-of-home consumption coefficients, selected years and food categories (LightGBM and TabPFN). Source: \texttt{model\_comparison\_national.csv}.}
\label{tab:national_coef_summary}
\begin{tabular}{llcccc}
\toprule
Model    & Year & Rice (daogu) & Wheat (xiaomai) & Pork (zhurou) & Vegetables (shucai) \\
\midrule
LightGBM & 2015 & 1.135 & 1.177 & 1.234 & 1.073 \\
         & 2020 & 1.128 & 1.153 & 1.181 & 1.069 \\
         & 2024 & 1.110 & 1.175 & 1.168 & 1.062 \\
\midrule
TabPFN   & 2015 & 1.139 & 1.213 & 1.313 & 1.124 \\
         & 2020 & 1.157 & 1.161 & 1.327 & 1.129 \\
         & 2024 & 1.233 & 1.149 & 1.376 & 1.187 \\
\bottomrule
\end{tabular}
\end{table}

\subsection{Imputation Model Selection}

Missing values in the out-of-home consumption ratio data are imputed using a supervised learning procedure. Model selection is based on hold-out performance on the imputation task. Metrics and ranks are reported in \texttt{data/imputation\_selection\_results.csv}.

The best-performing imputation model is FT-Transformer (\texttt{is\_best=True}): MAE = 0.1704, RMSE = 0.2978, with rank 1. ResNet Tabular ranks second (MAE 0.1793, RMSE 0.3054), followed by MLP, TabM, TabPFN, Random Forest, CatBoost, LightGBM, Linear, Lasso, XGBoost, and TabNet (worst MAE 0.2512). The negative $R^2$ values in that file reflect out-of-sample evaluation and indicate that the imputation task is difficult; MAE and RMSE are the primary criteria for selecting the FT-Transformer as the best imputation model. All downstream national and provincial results that use imputation therefore rely on the FT-Transformer-imputed ratios unless otherwise stated.

\subsection{Robustness: Main Versus No-Imputation (Robust) Specifications}

To assess sensitivity to imputation, we re-estimate national coefficients using the same prediction models but without imputing missing ratios (robust runs). The main results are in \texttt{results\_national\_<model>.csv}; the no-imputation counterparts are in \texttt{results\_national\_<model>\_robust.csv}.

For LightGBM, the main and robust coefficients are close in level and trend. In 2015, the main rice coefficient is 1.135 and the robust 1.141 (\texttt{results\_national\_lightgbm.csv} and \texttt{results\_national\_lightgbm\_robust.csv}, column \texttt{全国户外消费系数\_daogu}); in 2024, main is 1.110 and robust 1.140. Pork in 2024 is 1.168 (main) vs.\ 1.129 (robust). Omitting imputation thus slightly raises the rice coefficient and slightly lowers the pork coefficient for this model.

For TabPFN, the robust coefficients are generally somewhat higher for grains and some other categories. In 2015, rice is 1.139 (main) vs.\ 1.221 (robust); in 2024, 1.233 (main) vs.\ 1.255 (robust). Vegetables in 2024 are 1.187 (main) vs.\ 1.173 (robust). For FT-Transformer, 2015 rice is 1.206 (main) vs.\ 1.245 (robust) and 2024 rice 1.242 (main) vs.\ 1.263 (robust). The direction and size of the main--robust gap vary by model and category, but in no case do we find a qualitative reversal: out-of-home consumption continues to add to total consumption (coefficients remain above one), and the magnitudes remain in a comparable range. We conclude that the main findings are robust to excluding imputation, with the caveat that point estimates can shift by a few percentage points depending on the model and food category.

\subsection{Model Performance on the Prediction Task}

Prediction performance for the at-home consumption ratio (the training target) is summarized in \texttt{model\_accuracy\_summary.csv}. These metrics reflect internal predictive accuracy on the microdata evaluation design used in the project pipeline; they do not constitute external validation against independent macro-level measurements, which are unavailable for China at the required resolution. In the summary file, TabPFN attains the lowest MAE (0.0797), followed by TabM (0.0852), FT-Transformer (0.0865), and LightGBM (0.0872). RMSE rankings are similar, with FT-Transformer having the lowest RMSE (0.1432). Category-level metrics in \texttt{model\_accuracy\_detail.csv} show that performance varies by food category; for instance, FT-Transformer achieves the highest $R^2$ for rice (q\_daogu) among the models considered. Taken together with cross-model stability in the national and provincial outputs, these results motivate our emphasis on LightGBM, TabPFN, and FT-Transformer when presenting headline estimates.

\subsection{Uncertainty Quantification}

Bootstrap prediction intervals are produced by the project pipeline (\texttt{run\_bootstrap.py}): the prediction step is run repeatedly with different random seeds, and the 2.5\%--97.5\% percentiles across runs yield 95\% intervals for each province--year--category cell. Full bootstrap-based 95\% intervals for all models and categories are stored in the replication outputs \texttt{predictions\_<model>\_bootstrap.csv} (columns \texttt{*\_Mean}, \texttt{*\_Lower}, \texttt{*\_Upper}). These intervals confirm that the main qualitative conclusion---coefficients above one for all categories---is unchanged when uncertainty is accounted for, while also making transparent the residual model and sampling uncertainty for policy applications.

\subsection{Provincial Heterogeneity}

Provincial out-of-home consumption coefficients are estimated for all 12 models and stored in \texttt{results\_province\_*.csv}. Table~\ref{tab:provincial_heterogeneity} summarizes coefficient heterogeneity across selected provinces and years for the TabPFN model (source: \texttt{results\_province\_tabpfn.csv}, columns \texttt{q\_daogu\_coef}, \texttt{q\_zhurou\_coef}, \texttt{q\_shucai\_coef}). Province codes follow the project convention (GB/T 2260): 11 = Beijing, 13 = Hebei, 21 = Liaoning. Coefficients vary substantially by province: for example, in 2024 rice (daogu) ranges from about 1.13 (province 21) to 1.53 (province 11), and pork (zhurou) from about 1.30 to 1.38. Full provincial results for all models, years, and categories are available in the replication outputs \texttt{results\_province\_*.csv}.

\begin{table}[htbp]
\centering
\caption{Provincial out-of-home consumption coefficients (TabPFN), selected provinces, years and categories. Source: \texttt{results\_province\_tabpfn.csv}.}
\label{tab:provincial_heterogeneity}
\begin{tabular}{llcccc}
\toprule
Province & Year & Rice (daogu) & Wheat (xiaomai) & Pork (zhurou) & Vegetables (shucai) \\
\midrule
11 & 2020 & 1.428 & 1.160 & 1.279 & 1.272 \\
11 & 2024 & 1.533 & 1.208 & 1.377 & 1.397 \\
21 & 2020 & 1.090 & 1.049 & 1.246 & 1.110 \\
21 & 2024 & 1.129 & 1.051 & 1.300 & 1.149 \\
13 & 2020 & 1.197 & 1.142 & 1.326 & 1.149 \\
13 & 2024 & 1.281 & 1.143 & 1.375 & 1.216 \\
\bottomrule
\end{tabular}
\end{table}
